\newpage
\pagenumbering{roman}
\setcounter{page}{2}

\phantomsection
\addcontentsline{toc}{section}{Апстракт}

\begin{center}
    {\Large \textbf{Апстракт}}
\end{center}

\vspace{0.5cm}

Во овој дипломски труд обработувам практична и актуелна тема од
денешницата: следење и подобрување на здравата работна околина и
поддршка на здравјето на вработените во модерните работни простори.
Секојдневието на инженерите често вклучува долги часови работа во
затворени простории, при што малите промени во квалитетот на воздухот и
удобноста можат да предизвикаат замор, намалена концентрација и пад во
ефикасноста. Потребен е систем кој автоматски ги следи овие параметри и
овозможува навремена анализа, со што и навремена реакција за подобрување
на состојбата. Така во рамките на овој дипломски труд изработувам IoT
(Internet of Things) систем која прибира податоци од повеќе различни
извори: сензори за околината, Android паметни уреди како и надворешни
сервиси за временска прогноза и загадување. Системот обезбедува
автоматизирано следење, собирање, визуелизација како и обработка со
помош на вештачка интелигенција. Платформата е изработена врз
микро-сервисна архитектура, користејќи MQTT како главен протокол за
комуникација измеѓу уредите и платформата, взаемна автентикација со TLS,
\texttt{Redis Streams} за непречена проток на отчитувањата до \texttt{TimescaleDB} базата
на податоци, а врз сето тоа развиена е и веб контролна табла, сервис за
интелигентна анализа на собраните податоци со помош на вештачка
интелигенција, како и Android апликација што овозможува прибирање на
фитнес податоци. На овој начин, развиената платформа претставува
целокупен пристап за објективно мерење, анализа и интерпретација на
различните фактори кои влијаат врз продуктивноста и здравјето во
модерните работни простории.

\vspace{0.5cm}

\textbf{Клучни зборови:} \textit{IoT, MQTT, mTLS, микросервиси, TimescaleDB, вештачка интелигенција, работна околина, телеметрија.}
